\chapter[REVIEW OF RELATED LITERATURE]{Review of Related Literature}\label{ch:II}

\section{Technological Trends in Mapping Applications}\label{sec:2.1}

Location-based services (LBS) is an umbrella term for technologies that use a
device's location to provide services. Object positioning, navigation, and
mapping all count as LBS. Junglas and Watson~\cite{ref03} observed the
difference between LBS and non-LBS systems through an empirical study involving
navigation tasks, and it was found that individuals prefer LBS because it is
efficient and easy to use. However, these individuals also had privacy concerns:
they were willing to provide their location to network providers but hesitant to
give the same information to other entities such as third-party apps.

Another technology that is closely related to LBS is the Geographic Information
System (GIS). GIS is similar to LBS in that they both involve spatial and
location data, but one main difference lies in their intended purpose: GIS is
often used for mapping and analyzing geographic information, while LBS is more
user-oriented and focuses on navigation and on adding that same functionality to
everyday apps. GIS has made great progress over the past decades, but there are
still areas for improvement, such as geographic modeling enhancement and
human-oriented presentation of geographic information~\cite{ref04}.

LBS and GIS are far from the only technological trends. Kunhoth et
al.~\cite{ref05} provide an overview of various indoor positioning and
wayfinding methods from 2014 to 2020. Among these technologies are radio
frequency identification (RFID), ultra-wideband (UWB), and computer vision-based
indoor mapping.

\section{Commonly Used Navigation Apps}\label{sec:2.2}

To understand how navigation technologies work, we first take a look at Google
Maps, one of the most popular and feature-complete navigation apps. A report by
Mehta et al.~\cite{ref06} analyzes this app in terms of its navigation and
mapping. In terms of navigation, Google Maps uses the A* pathfinding algorithm.
For its mapping, Google uses specialized devices such as a 360° camera attached
to a car. This improves the accuracy of paths and provides 3D viewing through
Street View. To calculate the estimated time of arrival (ETA), Google Maps
collects data from its users on that particular route and approximates the
traffic speed. This means that the more users there are, the more accurate the
estimates will be. However, there are also ethical considerations here, and
transparency is important so that users know whether their personal information
is being used against them. Aside from Google Maps, an alternative mapping
solution is OpenStreetMap, a community-driven digital map where volunteers can
contribute to its development~\cite{ref07}. OpenStreetMap emphasizes
collaboration, organizing ``mapping parties'' that annotate data on various
geographical areas.

\section{Campus Navigation Solutions}\label{sec:2.3}

There exist several studies that involve indoor mapping for universities around
the world. Vanjani et al.~\cite{ref08} implemented a navigation app for Jhulelal
Institute of Technology (JIT) in Nagpur, India. The application was created to
help college students reach JIT College from various outside places
(inter-college), as well as offer directions to buildings and classrooms within
the university (intra-college). For the inter-college stage, they used GPS and
the Apple Maps API with a fixed destination towards JIT College. For the
intra-college stage, the paths to and from each building were hard-coded, and a
pre-made image was shown on the screen. This method works, but it is not ideal
since the number of possible paths is close to \emph{n}\textsuperscript{2}
(where \emph{n} is the number of buildings), so having a large number of
buildings increases the manual work quadratically.

There are also several navigation studies in the Philippines. Castillo et
al.~\cite{ref09} created a mobile navigation app for the Technological Institute
of the Philippines (T.I.P.) Quezon City campus, intended for new students and
visitors. First, the user is taken to a login screen. After creating an account
and logging in, the user can choose a source and a destination building, and the
app generates a path between them using a fuzzy logic algorithm. The advantage
of this feature over~\cite{ref08} is that the paths are automatically generated.
However, it may be prone to errors depending on the strength of the algorithm.
As an added feature, the user can also see other active users (shown as
pinpoints) within the same area using a channel selection algorithm through
Bluetooth.

Similar to~\cite{ref09}, Ramos et al.~\cite{ref10} created a locator map of
Saint Michael's College of Laguna using Android Studio. It uses the A* algorithm
for pathfinding. The layout of the campus map, including the buildings and the
walkable roads, was manually created. The app also has a search bar where the
user can search for a particular building, and it displays information about
that building, including the classes one needs to take in that building.
Overall, the app was rated ``Excellent'' based on 346 respondents.

A more gamified approach to campus navigation is found in the JUAN app of Don
Mariano Marcos Memorial State University~\cite{ref11}. The mobile campus map
incentivizes the user to go to different locations within the campus to collect
BeeCoins, which represent the number of points the user has. Through this
approach, the user can have fun while learning. The app used satellite view for
the map layout and GPS for the user's location.

\section{Indoor Navigation in Other Industries}\label{sec:2.4}

Aside from campus navigation, indoor mapping can also be effective in other
settings, such as finding stores in malls. A system proposed by Tulaskar et
al.~\cite{ref12} uses Wi-Fi fingerprinting to locate the user and Dijkstra's
algorithm for pathfinding. The application was developed using Android Studio
and Node.js, and it can display directions from the user's location to other
stores in the mall. However, it is unclear how multi-floor navigation
functionality can be achieved based on this implementation.

Another implementation by Saeliw et al.~\cite{ref13} combines augmented reality
(AR), geographic information systems (GIS), and GPS technologies to find stores.
To do this, a database was kept storing data on the floor level, latitude, and
longitude of each store. The user could then use their phone camera to locate
the store through an AR popup. Again, it is unclear whether the floor level of
the user was automatically determined or manually inputted.

Similar to~\cite{ref12}, Tan et al.~\cite{ref14} use Wi-Fi instead of GPS to
improve the accuracy of the user's indoor location. The mobile application used
an existing indoor map found in Google Maps and provided an interface for
locating stores inside Plaza Low Yat in Kuala Lumpur, Malaysia. On a scale of 1
to 5, users consistently rated it above 4 in terms of its usefulness, ease of
use, efficiency, and aesthetics.

\section{Challenges of Indoor Navigation}\label{sec:2.5}

While indoor navigation is certainly beneficial to all kinds of users, it does
not come without its challenges. A study by Zlatanova et al.~\cite{ref15}
proposes a set of indoor mapping and modeling (IMM) issues that are worth
researching. Among these problems are optimal navigation routing, emergency
preparedness, and legal and security issues.

Indoor navigation is believed to promote accessibility, especially for users
with mobility challenges. However, there is still a lack of actual testing with
these individuals. According to an analysis of 111 existing wayfinding and
navigation papers~\cite{ref16}, persons with mobility issues had no involvement
in the design and development of the navigation solutions in most papers. This
makes it hard to prove the effectiveness of the product in real-life contexts.

Another problem with indoor navigation is its lack of accessibility standards
and guidelines. Striegl et al.~\cite{ref17} report that in OpenStreetMap, on
average, only one building per city has an indoor map, and among those
buildings, the ones with accessibility information are even fewer.

One of the biggest challenges of indoor positioning, mapping, and navigation
(IPMN) is its lack of standardization. Over the decades, multiple researchers
and developers have tackled IPMN, but it is hard to undergo the process of
industrialization because there is no clear set of rules to be followed.
Gelernter and Maheshwari~\cite{ref18} also point out the lack of a standardized
file format for indoor map data, making it harder to have a one-size-fits-all
approach to speed up the indoor mapping process globally. To address this, Deng
et al.~\cite{ref19} proposed a set of requirements for IPMN:

\begin{enumerate}
\def\labelenumi{\arabic{enumi}.}
\item
  Indoor positioning can be classified into two types: requiring external
  equipment, or device-free. If using the former, the external equipment must be
  prepared in advance.
\item
  Indoor positioning must not be separated from map visualization, and map
  visualization should consider expressing symbols, color, and semantic
  information.
\item
  Indoor and outdoor positioning should be treated as separate modes in the
  application.
\item
  A consistency test is required for evaluating the accuracy and cost of
  positioning technologies.
\end{enumerate}
